\documentclass[11pt,a4paper]{article}

\usepackage[utf8]{inputenc}
\usepackage[T1]{fontenc}
\usepackage{geometry}
\usepackage{hyperref}
\usepackage{enumitem}
\usepackage{xcolor}
\usepackage{fancyhdr}
\usepackage{amsmath}
\usepackage{booktabs}
\usepackage{listings}

\geometry{margin=1in}
\hypersetup{
  colorlinks=true,
  linkcolor=blue,
  urlcolor=blue
}

\setlist{itemsep=0.35em, topsep=0.35em}

\pagestyle{fancy}
\fancyhf{}
\lhead{Object Oriented Programming (B.Tech CSE)}
\rhead{Unit 03: Diagnostic Assessment}
\cfoot{\thepage}
\setlength{\headheight}{14pt}

\lstset{
  language=Java,
  basicstyle=\ttfamily\small,
  keywordstyle=\color{blue},
  commentstyle=\color{gray},
  stringstyle=\color{teal},
  showstringspaces=false,
  columns=fullflexible,
  frame=single
}

% Marks per question (total = 100)
\newcommand{\EMarks}{\textbf{[2 marks]} }
\newcommand{\MMarks}{\textbf{[3 marks]} }
\newcommand{\HMarks}{\textbf{[5 marks]} }

\begin{document}

\begin{center}
  {\LARGE \textbf{Diagnostic Assessment -- Unit 03}}\\[0.25em]
  {\large Object Oriented Programming (CSEG1044)}\\[0.25em]
  \normalsize (Nested classes; exceptions + handlers; multithreading + synchronization; Java I/O streams)
\end{center}

\noindent
\textbf{Instructor:} Tofik Ali \hfill \textbf{Time Limit:} None\\
\textbf{Submission Deadline:} March 31, 2026\\
\textbf{Student Name:} \_\_\_\_\_\_\_\_\_\_\_\_\_\_\_\_\_ \hfill \textbf{Roll No.:} \_\_\_\_\_\_\_\_\\

\vspace{0.6em}
\hrule
\vspace{0.8em}

\textbf{Instructions}
\begin{itemize}[leftmargin=*]
  \item This assessment covers Unit-03 topics:
    \textbf{Nested classes}, \textbf{exception handling/handlers}, \textbf{multithreading + synchronization},
    and \textbf{Java I/O streams}.
  \item \textbf{Total marks: 100.} Easy: 10 $\times$ 2 = 20, Medium: 10 $\times$ 3 = 30, Hard: 10 $\times$ 5 = 50.
  \item \textbf{No time limit.} Submit on or before \textbf{March 31, 2026}.
  \item \textbf{Important (70\% rule): To have your assignment marks counted, you must score at least 70\% (70/100 or more) in this assessment.}
    If you score below 70\%, you may reattempt and resubmit. Your latest submission will be considered.
  \item Answer \textbf{all} questions. For code questions, snippets are acceptable unless a full program is requested.
  \item For MCQs, follow the instruction in the question (some are \textbf{select all that apply}).
\end{itemize}

\vspace{0.6em}

\section*{Easy (10 Questions)}
\begin{enumerate}[label=\textbf{E\arabic*.}, leftmargin=*]
  \item \EMarks \textbf{[MCQ -- Select all that apply]} Which are types of nested classes in Java?
    \begin{enumerate}[label=(\Alph*), leftmargin=*]
      \item Static nested class
      \item Inner (non-static) class
      \item Local class
      \item Anonymous class
      \item Global class
    \end{enumerate}

  \item \EMarks \textbf{[Subjective]} Does an inner class require an outer object? Explain in 2--3 lines.

  \item \EMarks \textbf{[MCQ]} Which of the following is a \textbf{checked} exception?
    \begin{enumerate}[label=(\Alph*), leftmargin=*]
      \item \texttt{IOException}
      \item \texttt{NullPointerException}
      \item \texttt{ArithmeticException}
      \item \texttt{StackOverflowError}
    \end{enumerate}

  \item \EMarks \textbf{[MCQ]} Which block always executes (if the program doesn't terminate abruptly)?
    \begin{enumerate}[label=(\Alph*), leftmargin=*]
      \item \texttt{try}
      \item \texttt{catch}
      \item \texttt{finally}
      \item \texttt{throw}
    \end{enumerate}

  \item \EMarks \textbf{[Subjective]} Difference between \texttt{throw} and \texttt{throws} (3--4 lines).

  \item \EMarks \textbf{[MCQ]} Which method starts a new thread of execution?
    \begin{enumerate}[label=(\Alph*), leftmargin=*]
      \item \texttt{run()}
      \item \texttt{start()}
      \item \texttt{sleep()}
      \item \texttt{wait()}
    \end{enumerate}

  \item \EMarks \textbf{[MCQ -- Select all that apply]} Valid ways to create a thread in Java are:
    \begin{enumerate}[label=(\Alph*), leftmargin=*]
      \item Extend \texttt{Thread}
      \item Implement \texttt{Runnable}
      \item Call \texttt{run()} directly to start a new OS thread
      \item Change the \texttt{main} method signature
    \end{enumerate}

  \item \EMarks \textbf{[MCQ]} What is the main purpose of \texttt{synchronized}?
    \begin{enumerate}[label=(\Alph*), leftmargin=*]
      \item Makes code execute faster
      \item Ensures mutual exclusion for shared data
      \item Automatically creates threads
      \item Reads data from a file
    \end{enumerate}

  \item \EMarks \textbf{[MCQ]} \texttt{FileReader} is mainly used for:
    \begin{enumerate}[label=(\Alph*), leftmargin=*]
      \item Byte/binary data
      \item Character/text data
      \item Network packets
      \item Graphics rendering
    \end{enumerate}

  \item \EMarks \textbf{[MCQ]} What does \texttt{try-with-resources} guarantee?
    \begin{enumerate}[label=(\Alph*), leftmargin=*]
      \item No exceptions will occur inside the \texttt{try} block
      \item The resource is closed automatically
      \item I/O becomes faster automatically
      \item Threads stop automatically on error
    \end{enumerate}
\end{enumerate}

\section*{Medium (10 Questions)}
\begin{enumerate}[label=\textbf{M\arabic*.}, leftmargin=*]
  \item \MMarks \textbf{[Code]} Write code to instantiate a static nested class \texttt{Outer.Nested}.

  \item \MMarks \textbf{[Code]} Write code to instantiate an inner class \texttt{Outer.Inner}.

  \item \MMarks \textbf{[Code]} Write a \texttt{try/catch/finally} block that handles divide-by-zero and always prints \texttt{Done}.

  \item \MMarks \textbf{[Code]} Write a method \texttt{readFirstLine(String path) throws IOException} and show how to call it using \texttt{try/catch}.

  \item \MMarks \textbf{[Code]} Create a thread using \texttt{Runnable} that prints 1--3 and start it.

  \item \MMarks \textbf{[Concept]} What is a race condition? Explain in 4--6 lines.

  \item \MMarks \textbf{[Code]} Write a thread-safe counter increment method using \texttt{synchronized}.

  \item \MMarks \textbf{[Code]} Write a \texttt{try-with-resources} snippet that writes ``Hello'' to \texttt{msg.txt} using \texttt{FileWriter}.

  \item \MMarks \textbf{[Code]} Write a snippet that reads \texttt{msg.txt} using \texttt{FileReader} and prints its content.

  \item \MMarks \textbf{[Concept]} What are thread priorities? Do they guarantee execution order? (4--6 lines)
\end{enumerate}

\section*{Hard (10 Questions)}
\begin{enumerate}[label=\textbf{H\arabic*.}, leftmargin=*]
  \item \HMarks \textbf{[Program]} Write a complete program that demonstrates try/catch/finally execution sequence for:
    (a) success case, (b) exception case. Include output.

  \item \HMarks \textbf{[Program]} Write a program with two threads incrementing a shared counter (e.g., 100000 times each).
    First run \textbf{without} synchronization, then fix using \texttt{synchronized}. Print final values.

  \item \HMarks \textbf{[Program]} Ticket booking: simulate \texttt{availableSeats = 5} shared by 10 threads (each books 1).
    Use synchronization to prevent overbooking and print final result.

  \item \HMarks \textbf{[Program]} Write a program that prints thread states using \texttt{getState()} to show:
    \textbf{NEW}, \textbf{RUNNABLE}, \textbf{TIMED\_WAITING} (via \texttt{sleep}), and \textbf{TERMINATED}.

  \item \HMarks \textbf{[Program]} Copy a text file using \texttt{FileReader}/\texttt{FileWriter} and a character buffer.
    Use \texttt{try-with-resources} and handle \texttt{IOException}.

  \item \HMarks \textbf{[Program]} Read \texttt{numbers.txt} (one value per line, some invalid).
    Sum valid integers, report invalid lines, and write sum to \texttt{sum.txt}.

  \item \HMarks \textbf{[Program]} Create an outer class with a static nested class \texttt{Task implements Runnable}.
    Start two tasks concurrently and print results.

  \item \HMarks \textbf{[Program]} Show exception propagation with I/O using a 3-method call chain:
    the deepest method declares \texttt{throws IOException} and \texttt{main} handles it.

  \item \HMarks \textbf{[Explain]} Differentiate byte streams vs character streams and justify why
    \texttt{FileReader}/\texttt{FileWriter} are appropriate for text files (8--12 lines).

  \item \HMarks \textbf{[Program]} Write code showing that \texttt{finally} runs even when \texttt{return} is executed in try.
    Include output.
\end{enumerate}

\vfill
\noindent\textit{End of Assessment}

\end{document}
